\section{The source harmonic space and its column orbit}
\label{sec:sr-source-harmonic-orbit}

This section gives the algebraic producer used by the short-family
construction.  In particular, the assertion that a column orbit spans
the required harmonic space is proved here, with the source's mixed
action and its integer equations.

Let \(K\) be a field of characteristic zero, \(T\) a finite set, and
\(p,q,n\ge0\) with \(p+q\le n\).  Work in the polynomial ring
\[
 {\cal P}_K=K[x_{t,a,k},y_{t,b,k}:
     t\in T,\ 1\le a\le p,\ 1\le b\le q,\ 1\le k\le n].
\]
For nonnegative arrays \(c_{t,i}\), \(c'_{t,j}\), set
\[
 a_{t,i}=\sum_{h=i}^{p}c_{t,h},\qquad
 b_{t,j}=\sum_{h=j}^{q}c'_{t,h},\qquad
 D=\sum_t\left(\sum_i i c_{t,i}+\sum_j j c'_{t,j}\right).
 \tag{HO1}\label{eq:sr-source-harmonic-degree}
\]
Thus \(a_t,b_t\) are antitone and \(D=\sum_{t,i}a_{t,i}
+\sum_{t,j}b_{t,j}\).  Define \({\cal H}_K(c,c')\) by the following
linear equations:
\[
 \Delta_{t,a,b}P=\sum_k
           \partial_{x_{t,a,k}}\partial_{y_{t,b,k}}P=0;
 \quad
 \deg_{x_{t,a,*}}P=a_{t,a},\quad
 \deg_{y_{t,b,*}}P=b_{t,b};
 \tag{HO2}\label{eq:sr-source-row-highest}
\]
the degree assertions mean that every supported monomial has the
indicated row degree.  Also require invariance under
\(x_{t,b,*}\mapsto x_{t,b,*}+s x_{t,a,*}\) for \(a<b\), and the
same substitutions in the \(y\) block, initially for every
\(s\in\mathbb Z\).  A polynomial identity in \(s\), over a
characteristic-zero field, makes this equivalent to invariance for
every \(s\in K\).  This is the source
\(\mathrm{rowHighest}(K,c,c')\).

For vectors \(\epsilon_{t,i},\phi_{t,j}\in K^n\), put
\[
 X_{t,a}(\epsilon)=\sum_k\epsilon_kx_{t,a,k},\qquad
 Y_{t,b}(\phi)=\sum_k\phi_ky_{t,b,k},
\]
and define the leading-minor product
\[
 {\cal M}_{c,c'}(\epsilon,\phi)=
 \prod_t\left\{
   \prod_{i=1}^{p}
       \det(X_{t,a}(\epsilon_{t,h}))_{1\le a,h\le i}^{c_{t,i}}
   \prod_{j=1}^{q}
       \det(Y_{t,b}(\phi_{t,h}))_{1\le b,h\le j}^{c'_{t,j}}
          \right\}.
 \tag{HO3}\label{eq:sr-source-minor-product}
\]
The source column action of \(g=(g_t)\in\prod_t\mathrm{GL}_n(K)\)
sends coefficient covectors by \(\epsilon\mapsto\epsilon g_t^{-1}\)
and coefficient vectors by \(\phi\mapsto g_t\phi\).
Let \(P_0={\cal M}_{c,c'}((e_i)_{i\le p},
(e_{n+1-j})_{j\le q})\), with these same coordinate vectors in
each \(t\) block.  In particular \(P_0\) has integer coefficients
and is nonzero.

\begin{srproposition}\label{prop:sr-source-harmonic-column-orbit}
For all the displayed choices,
\[
 {\cal H}_K(c,c')
       =\operatorname{span}_K\{gP_0:g\in\textstyle\prod_t
                                         \mathrm{GL}_n(K)\}.
 \tag{HO4}\label{eq:sr-source-harmonic-column-orbit}
\]
It is also the span of the iterated column operators applied to
\(P_0\), and it is spanned over \(K\) by the images of
\({\cal H}_{\mathbb Z}(c,c')\).  Empty blocks and \(n=0\) are
allowed; if \(n=0\), the inequality \(p+q\le n\) forces
\(p=q=0\), and the assertion is the statement that constants
span \(K\).
\end{srproposition}

\begin{proof}
We supply the source's finite-dimensional highest-vector argument.

Write, for fixed \(t\),
\[
 E_{kl}=\sum_a x_{t,a,k}\partial_{x_{t,a,l}}
          -\sum_b y_{t,b,l}\partial_{y_{t,b,k}}.
 \tag{HO5}\label{eq:sr-source-mixed-column-operator}
\]
The identity
\[
 [z_u\partial_{z_v},z_w\partial_{z_s}]
       =\delta_{v,w}z_u\partial_{z_s}
                      -\delta_{s,u}z_w\partial_{z_v}
\]
gives \([E_{ij},E_{kl}]=\delta_{j,k}E_{il}
-\delta_{l,i}E_{kj}\).  Operators for distinct \(t\)'s commute.
The column action preserves all equations in (HO2), by the
chain rule and \(g_t^{-1}g_t=1\), and commutes with the two row
unipotent actions.  These facts also follow by checking each
variable and then applying the product rule.

The orbit span is stable under every \(E_{kl}\).
For \(k\ne l\), the substitution for the transvection
\(g_t=1-s\,e_{lk}\) has first coefficient \(E_{kl}P\);
coefficient extraction from this vector-valued polynomial in
\(s\) preserves any subspace containing all its values.
For \(k=l\), take the diagonal matrix with \(k\)-th entry
\(u^{-1}\) and all other diagonal entries \(1\).
Under the displayed mixed action it acts on a monomial \(m\) by
\(u^{w_{t,k}(m)}\), where
\[
 w_{t,k}(m)=\sum_a m(x_{t,a,k})-\sum_bm(y_{t,b,k})\in\mathbb Z.
\]
Multiplying the finite Laurent sum by a power of \(u\), and
using polynomial interpolation on \(K^\times\), extracts its
weight components and their integer-weighted sum.
That sum is \(E_{kk}P\).
The base polynomial and every column translate lie in
\({\cal H}_K\): a contraction differentiating one \(x\) and
one \(y\) factor has coefficient \(e_i\cdot e_{n+1-j}=0\);
row degrees and row-unipotent invariance follow directly
from determinant multilinearity.
Thus every word in the \(E_{kl}\) applied to \(P_0\) is both
in the orbit span and in \({\cal H}_K\).

We next prove that every nonzero subspace
\(W\subseteq{\cal H}_K\) stable under all the \(E_{kl}\)
contains \(P_0\).  The ambient degree is \(D\), so \(W\) is
finite dimensional.  The diagonal \(E_{t,kk}\) commute and
are diagonal on monomials with integer eigenvalues.
Their finitely many eigenprojections are polynomials in
these operators; they therefore preserve \(W\).
There is a nonzero joint eigenvector \(v\in W\), with
integer weights \(\kappa_{t,k}\), minimizing
\(\sum_{t,k}k\kappa_{t,k}\) among the weights occurring in
\(W\).  For \(k<l\), a nonzero \(E_{t,kl}v\) would have this
height decreased by \(l-k\).  Consequently
\(E_{t,kl}v=0\) for \(k<l\).

For completeness, these highest weights are antitone.
Take \(E=E_{t,ij}\), \(F=E_{t,ji}\),
\(H=E_{t,ii}-E_{t,jj}\), \(i<j\), and
\(\ell=\kappa_{t,i}-\kappa_{t,j}\).
The commutators give
\[
 H F^kv=(\ell-2k)F^kv,\qquad
 E F^{k+1}v=(k+1)(\ell-k)F^kv.
\]
If \(\ell<0\), induction makes every \(F^kv\) nonzero;
their distinct \(H\)-weights make them independent, contrary
to finite dimensionality.  Thus \(\ell\ge0\).

We show directly that \(v\) is a scalar multiple of \(P_0\).
The following row-degree bound is the first ingredient.
If a polynomial of row degrees \(a_1,\ldots,a_p\) is invariant
under adding an earlier row to a later row, then any supported
monomial \(m\), and any set \(S\) of \(s\) columns, satisfy
\[
 \sum_{a,k\in S}m(x_{a,k})\le\sum_{a\le\min(s,p)}a_a.
 \tag{HO6}\label{eq:sr-source-column-degree-dominance}
\]
For \(s\ge p\) this follows from total row degrees.
For \(s<p\), on the locus where the first \(s\) rows in
columns \(S\) have nonzero determinant, add their linear
combinations to later rows to clear those columns.
These operations commute with scaling the columns \(S\)
by a variable \(z\).  Invariance makes the coefficient of
any power of \(z\) greater than the right side vanish on
that locus.  Let \(R\) be such a coefficient polynomial and
\(\Delta\) this nonzero determinant.  At every point
\(R\Delta=0\), including points with \(\Delta=0\);
polynomial evaluation over the infinite field gives
\(R\Delta=0\), hence \(R=0\).
This proves (HO6), including its assertion for each supported
monomial.

Fix \(t\), suppress it in the notation, and set
\[
 \lambda=(a_1,\ldots,a_p,0,\ldots,0,-b_q,\ldots,-b_1).
\]
The diagonal eigenvalue equations imply that every monomial
of \(v\) has column weight \(\kappa\).
Apply (HO6) to the first \(i\) columns of the \(x\) block
and the last \(n-i\) columns of the \(y\) block.  Since the
total \(y\)-degree is \(\sum b_j\), this gives
\[
  \sum_{k\le i}\kappa_k\le\sum_{k\le i}\lambda_k,
       \qquad \sum_k\kappa_k=\sum_k\lambda_k.
 \tag{HO7}\label{eq:sr-source-prefix-weight-bound}
\]

The second ingredient is the explicit Casimir identity.
Let \(R^x_{ab}=\sum_kx_{a,k}\partial_{x_{b,k}}\),
and similarly \(R^y_{ab}\).
The product rule and commuting partial derivatives give
\[
\begin{split}
 \sum_{k,l}E_{kl}E_{lk}
 ={}&\sum_{a,b}R^x_{ab}R^x_{ba}
           +(n-p)\sum_aR^x_{aa}\\
 &+\sum_{a,b}R^y_{ab}R^y_{ba}
           +(n-q)\sum_bR^y_{bb}
 -2\sum_{a,b}\left(\sum_kx_{a,k}y_{b,k}\right)\Delta_{a,b}.
\end{split}
 \tag{HO8}\label{eq:sr-source-capelli-casimir}
\]
To verify the unmixed term, expand both sums of compositions.
Their second-derivative terms are identical after exchanging
indices, while differentiating a coefficient contributes
\(n\sum x\partial_x\) on the column side and
\(p\sum x\partial_x\) on the row side.
The two mixed terms each give
\(\sum_{a,b}(\sum_kx_{a,k}y_{b,k})\Delta_{a,b}\).
This proves (HO8) with the displayed sign.

If \(C(\mu)=\sum_k\mu_k(\mu_k+n+1-2k)\), the highest-column
equations give Casimir scalar \(C(\kappa)\).
Indeed, for \(i<j\),
\(E_{ij}E_{ji}v=[E_{ij},E_{ji}]v
=(\kappa_i-\kappa_j)v\), whereas \(E_{ji}E_{ij}v=0\).
The same calculation for the two row actions, using their
degrees and row-highest equations, and then (HO8) with
\(\Delta_{a,b}v=0\), gives scalar
\[
 \sum_i a_i(a_i+n+1-2i)+
 \sum_j b_j(b_j+n+1-2j)=C(\lambda).
\]
Thus \(C(\kappa)=C(\lambda)\).
This equality and (HO7) force \(\kappa=\lambda\):
put \(S_i=\sum_{k\le i}(\lambda_k-\kappa_k)\).
Then \(S_i\ge0\), \(S_n=0\), and summation by parts gives
\[
 C(\lambda)-C(\kappa)
  =\sum_{i<n}S_i
    \bigl(\lambda_i+\kappa_i-\lambda_{i+1}-\kappa_{i+1}+2\bigr).
\]
Both weights are antitone, so each parenthesis is at least
\(2\).  Equality forces all \(S_i=0\).

For a monomial of \(v\), the first \(p\) column weights
now sum to its total \(x\)-degree.  Nonnegative exponents
therefore force all \(x\)'s outside the first \(p\) columns
to have exponent zero.  Similarly all \(y\)'s outside the
last \(q\) columns have exponent zero.
The highest-column equations reduce to ordinary
column-unipotent invariance on the \(p\)-square \(x\) block
and on the \(q\)-square \(y\) block with its columns
reversed.  In characteristic zero, annihilation by the
corresponding row-addition derivation is equivalent to
invariance: for the generic addition \(A(s)\),
\((A(s)P)'=A(s)(DP)\), so \(DP=0\) makes \(A(s)P\)
constant.  Conversely, differentiation at zero gives
\(DP=0\).

A polynomial bi-invariant under earlier-row and
earlier-column additions, and with row degrees
\(a_i=\sum_{h\ge i}c_h\), is
\[
 Q\prod_i\Delta_i^{c_i},
 \tag{HO9}\label{eq:sr-source-leading-minor-factor}
\]
where \(Q\) is independent of that square block.
Here is the source's elimination proof.
If all leading minors are nonzero, successive row and
column clearings diagonalize the block while preserving
every bi-invariant function and every leading minor.
At a diagonal block \(\mathrm{diag}(d_i)\), row degrees
give the value at the identity block times
\(\prod_i d_i^{a_i}=\prod_i(d_1\cdots d_i)^{c_i}\).
Take \(Q\) to be the polynomial obtained by replacing
this block by the identity.  To extend the resulting
equality to all blocks, multiply its difference by the
product of the leading minors and use polynomial
evaluation over \(K\) as in the proof of (HO6).
That product is nonzero, so cancellation proves (HO9).
The argument also works for a square block of size zero.

Apply (HO9) to each of the disjoint \(x\) and reversed
\(y\) blocks, for every \(t\).  There are no remaining
variables in \(v\).
Since the quotient for each block is independent of that
block and its minor product has value \(1\) at the identity,
replacing that block gives the identity
\(v(x)=v(x\text{ with that block replaced by }1)D_{\rm block}(x)\).
Replacing all these blocks successively by identity
matrices shows
\(v=v(x_0)P_0\), where \(P_0(x_0)=1\).
This proves the claimed uniqueness of the highest vector.
Since \(v\ne0\), its scalar is nonzero, and \(P_0\in W\).

We finish first over \(\mathbb C\).
Use the Fock inner product
\[
  \langle P,Q\rangle_F=\sum_m m!\,\overline{P_m}Q_m,
       \qquad m!=\prod_v m(v)! .
\]
It is positive definite, and the elementary
multiplication/derivation adjoint identity gives
\(E_{kl}^*=E_{lk}\).
Let \(U\) be the span of all column words applied to
\(P_0\), and let \(W={\cal H}_{\mathbb C}\cap U^\perp\).
The adjoint identity makes \(W\) stable under every
\(E_{kl}\), because \(U\) is stable under its adjoint.
If \(W\ne0\), the preceding argument puts \(P_0\) in
\(W\), although \(P_0\in U\), contradicting positivity.
Thus \(W=0\).  Orthogonal projection onto the finite
dimensional \(U\) now gives \({\cal H}_{\mathbb C}=U\).

Finally the equations defining \({\cal H}\), and all the
column words in \(P_0\), have integer coefficients.
For an extension of fields \(F\subseteq L\), applying
any \(F\)-linear functional \(L\to F\) coefficientwise
commutes with these equations.  A basis of \(L/F\)
expands any individual polynomial into a finite sum
of such coefficient polynomials.
Consequently
\(L\cdot{\cal H}_F={\cal H}_L\).
An \(F\)-linear functional \(L\to F\) taking \(1\) to
\(1\) also descends any finite spanning relation for
a polynomial originally over \(F\).
Apply this first to \(\mathbb Q\subset\mathbb C\), and
then to \(\mathbb Q\subset K\).  It proves the same
column-word spanning assertion over \(K\).
Those words have integer coefficients, yielding
the asserted integral spanning as well.
Since the words lie in the column orbit span and
that span lies in \({\cal H}_K\), (HO4) follows.
\end{proof}

We also record the precise orbit parametrization needed below.
The orbit of \(P_0\) consists of the polynomials (HO3)
with, for each \(t\), independent covectors
\(\epsilon_{t,1},\ldots,\epsilon_{t,p}\) and independent
vectors \(\phi_{t,1},\ldots,\phi_{t,q}\) satisfying
\(\epsilon_{t,a}\cdot\phi_{t,b}=0\).
To construct the required \(g_t\), view the
\(\epsilon\)'s as a surjection \(K^n\to K^p\).
The independent \(\phi\)'s lie in its kernel.
Extend them to a kernel basis and choose lifts of the
coordinate basis of \(K^p\); the resulting adapted basis
gives \(g_t\) with the required covector and vector
actions.

Arbitrary, possibly dependent, perpendicular arrays also
give a scalar multiple of an orbit member.
Indeed, since \(n-q\ge p\), extend the successive
covector spans to an independent flag of length \(p\)
in the annihilator of the original \(\phi\)'s.
Then, since \(n-p\ge q\), extend the successive original
vector spans to an independent flag in the annihilator
of these covectors.
For each leading determinant, the original leading
vectors lie in the corresponding new flag space;
their exterior product is a scalar multiple of its
basis wedge.  Multiplying those scalars, with the
exponents \(c,c'\), proves the assertion, including
a zero scalar when a required leading determinant
vanishes.  This is the source's flag-extension step,
so no independence premise will be imposed on the
finite grid.

\paragraph{Frozen-source correspondence.}
All paths below have prefix
\texttt{src/Development/AbelianVarietyTheory/}.
Definitions:
\texttt{Definitions/ThetaMetrics/IsWeightedHomogeneousOfMemOrbitSpan.lean}
(\texttt{c1877771faeabf2142790b50b81479d6d456888f});
\texttt{Definitions/ThetaValues/ExistsMemRowHighestTotalDegreeEqTwo.lean}
(\texttt{217fbb7c67c140ae738bd4cdf20ca1d889b544a0}).
Proofs under \texttt{Proofs/CuspFormEvaluations/}:
\texttt{SumSmulMemOfSumZpowSmulMem.lean}
(\texttt{3e061221e4036b034ee0f711d62a6eeac44c0f43});
\texttt{CommuteRowOpRightMixedOp.lean}
(\texttt{e5177dd3860b99ec938dcae8ed7e3bfc0fb6b883});
\texttt{ExistsJointIntEigenvectorOfStable.lean}
(\texttt{19583793c47ca3dc64c71ea7e82a1d1722e7a6e4});
\texttt{Sl2NonnegWeightOfHighestWeightVector.lean}
(\texttt{b7b5a34fe292d0bf1a4914f15a0b81ac91674edf});
\texttt{EqZeroOfPartialSumsNonneg.lean}
(\texttt{a0a1d943821745f91e4848a82fa606134a049104});
\texttt{ColumnDegreeLe.lean}
(\texttt{6228cdfba74d695455dae9204a0af5f24ceaf21e});
\texttt{RowAddGenericRowAddTerm.lean}
(\texttt{a330cf3dec19dd7558c4484d7396d47fcc32baac});
\texttt{ProdLeadMinorDiagonal.lean}
(\texttt{2efe5dbb526ffa538caf14b41c1fe4726237cf01});
\texttt{ExistsEqMulProdMinors.lean}
(\texttt{aa0719b052a50fa7ca15c2f4d00278833ca702eb});
\texttt{RowHighestLeSpanColWordComplex.lean}
(\texttt{6723d08a3376d7658011e5f319b0ee89cf7ef040}).
The final flag and orbit parametrization is
\texttt{Proofs/CovolumeBounds/ExistsMinorProdEqSmulColActBasePoly.lean}
(\texttt{2ad1227e78855888f65ff5874a3bb24bc58f6ff4}).
The declarations used are
\srcid{rowHighest_eq_span_colWord},
\srcid{rowHighest_eq_orbitSpan},
\srcid{span_map_rowHighest_int},
\srcid{casimir_value_strict}, and
\srcid{exists_minorProd_eq_smul_colAct_basePoly}.
